\chapter{Závěr}
\label{chap:conclusion}

Cílem této diplomové práce bylo provést komplexní marketingový výzkum zaměřený na posouzení potenciálu uvedení nové značky mraženého ovoce v~čokoládě Berrie na český trh. Na základě provedené analýzy ověřených sekundárních dat, doplněné o~odhadované spotřebitelské ukazatele odvozené z~oborových benchmarků, lze formulovat následující předběžné závěry. Definitivní závěry budou formulovány po dokončení primárního výzkumu.

% ============================================================================
\section{Hlavní zjištění}

\textbf{1. Tržní prostředí je příznivé pro vstup nového produktu.} Trh mražených zpracovaných ovoce a~zeleniny ve východní Evropě dosáhl v~roce 2024 hodnoty 2,53~mld.~USD s~meziročním růstem 7,4\,\% \parencite{frozenfoodeurope_2025}. Výdaje na osobu ve východní Evropě (8,10~USD) dosahují pouhé třetiny úrovně západní Evropy (24,00~USD), což signalizuje významný růstový potenciál. Spotřeba ovoce v~ČR se stabilně drží na úrovni 85--87~kg na osobu za rok \parencite{czso_spotreba_2024} a~spotřeba čokolády činí přibližně 6,5~kg na osobu za rok, což potvrzuje stabilní zájem českého spotřebitele o~obě klíčové složky produktu Berrie.

\textbf{2. Na českém trhu existuje jednoznačná mezera.} Segment mraženého ovoce v~prémiové čokoládě je v~současnosti obsluhován výhradně jedním zahraničním dodavatelem — argentinskou značkou Franui. Lokální alternativa neexistuje. Distribuce Franui je navíc omezena převážně na platformu Rohlík.cz, což vytváří prostor pro produkty s~širší distribucí.

\textbf{3. Cenová analýza potvrzuje prostor pro konkurenční vstup.} Výsledky PSM analýzy ukazují optimální cenový bod v~rozmezí 111--112~Kč za 150g balení, zatímco Franui se prodává za 159,99~Kč. Navrhovaná cena Berrie (119--130~Kč) spadá do přijatelného cenového rozmezí (92--131~Kč) a~nabízí spotřebiteli úsporu 19--26\,\% oproti konkurenci.

\textbf{4. Cílové segmenty vyžadují diferenciovanou komunikaci.} Dva identifikované cílové segmenty — generace~Z/mileniálové a~matky s~dětmi — se výrazně liší v~motivacích k~nákupu. Zatímco první segment reaguje primárně na vizuální atraktivitu a~sociální doporučení, druhý segment klade důraz na kvalitu surovin a~zdravotní aspekty.

% ============================================================================
\section{Doporučení pro praxi}

Na základě provedené analýzy se doporučuje:

\begin{enumerate}
  \item \textbf{Vstoupit na trh} s~produktem Berrie ve dvou vlajkových variantách (borůvky v~bílé čokoládě, jahody v~hořké čokoládě) s~cenovou hladinou 119--130~Kč za 150g balení.
  \item \textbf{Zvolit třífázovou distribuční strategii} počínající exkluzivním uvedením na Rohlík.cz (fáze~1), následovanou rozšířením do řetězců Albert a~Tesco (fáze~2) a~plným pokrytím trhu (fáze~3).
  \item \textbf{Investovat do influencer marketingu} na platformách TikTok a~Instagram s~důrazem na vizuálně atraktivní obsah, přičemž rozpočet pro první rok by měl činit přibližně 1,15~mil.~Kč.
  \item \textbf{Budovat značku na pilířích lokálního původu} (české ovoce, případně v~\gls{bio} kvalitě) a~technologické kvality (\gls{iqf} mražení).
\end{enumerate}

% ============================================================================
\section{Směry dalšího výzkumu}

Tato práce představuje první fázi výzkumného projektu. Pro kompletní validaci zjištění a~informované rozhodnutí o~vstupu na trh je nezbytné realizovat plánovaný primární výzkum:

\begin{itemize}
  \item \textbf{Skupinové diskuse (focus groups)} se dvěma cílovými segmenty pro kvalitativní validaci spotřebitelských preferencí a~bariér.
  \item \textbf{Kvantitativní dotazníkové šetření (\gls{cawi})} na vzorku 400 respondentů pro statisticky robustní verifikaci výsledků PSM analýzy a~měření nákupního záměru.
  \item \textbf{Senzorickou analýzu (\gls{clt})} pro ověření chuťových preferencí a~srovnání senzorického profilu Berrie s~konkurenčním Franui.
\end{itemize}

Realizace těchto fází je plánována na období červen--srpen 2026 a~jejich výsledky budou integrovány do finální verze této práce.

% ============================================================================
\section{Přínos práce}

Diplomová práce přináší několik praktických přínosů:
\begin{itemize}
  \item komplexní analýzu tržního potenciálu pro nový produktový segment na českém trhu,
  \item aplikaci metody Price Sensitivity Meter v~kontextu českého potravinářského trhu,
  \item návrh kompletní marketingové strategie (produkt, cena, distribuce, komunikace) podloženou daty,
  \item identifikaci klíčových faktorů úspěchu a~rizik pro vstup nového produktu do kategorie prémiových mražených pochutin.
\end{itemize}

Z~teoretického hlediska práce přispívá k~poznání spotřebitelského chování v~dynamicky rostoucím segmentu mražených ovoce v~ČR a~demonstruje využití smíšené výzkumné metodologie v~kontextu marketingového výzkumu pro nový produkt.
